\documentclass[12pt, a4paper]{article}
\usepackage[utf8]{inputenc}
\usepackage[T2A]{fontenc}
\usepackage[russian]{babel}
\usepackage{geometry}
\usepackage{booktabs}
\usepackage{array}
\usepackage{xcolor}
\usepackage{hyperref}
\usepackage{amsmath}
\usepackage{amssymb}
\usepackage{enumitem}
\usepackage{fancyhdr}
\usepackage{titlesec}
\usepackage{mdframed}
\usepackage{listings}
\setlength{\headheight}{13.6pt}

\geometry{top=2.5cm, bottom=2.5cm, left=3cm, right=2cm}

% --- Основные Цвета ---
\definecolor{primary}{RGB}{25, 118, 210}
\definecolor{success}{RGB}{46, 125, 50}
\definecolor{warning}{RGB}{230, 81, 0}
\definecolor{danger}{RGB}{183, 28, 28}
\definecolor{lightgray}{RGB}{245, 245, 245}
\definecolor{codegray}{RGB}{240, 240, 240}
\definecolor{darkgray}{RGB}{66, 66, 66}

% --- Фоновые цвета для блоков ---
\definecolor{bgsuccess}{RGB}{232, 245, 233}
\definecolor{bgwarning}{RGB}{255, 243, 224}
\definecolor{bgdanger}{RGB}{255, 235, 238}

\hypersetup{colorlinks=true, linkcolor=primary, urlcolor=primary}

% --- Заголовки секций ---
\titleformat{\section}{\Large\bfseries\color{primary}}{\thesection.}{0.5em}{}
[\vspace{-0.5em}\rule{\linewidth}{0.8pt}\vspace{0.3em}]
\titleformat{\subsection}{\large\bfseries\color{darkgray}}{\thesubsection.}{0.5em}{}
\titleformat{\subsubsection}{\normalsize\bfseries}{\thesubsubsection.}{0.5em}{}

\pagestyle{fancy}
\fancyhf{}
\fancyhead[L]{\small\color{darkgray} Отчёт по хакатону: LTV-предсказание}
\fancyhead[R]{\small\color{darkgray} Егоров Кирилл}
\fancyfoot[C]{\thepage}
\renewcommand{\headrulewidth}{0.4pt}

% --- Цветные блоков ---
\newmdenv[backgroundcolor=lightgray, linecolor=primary, linewidth=2pt,
leftline=true, rightline=false, topline=false, bottomline=false,
innerleftmargin=10pt, innertopmargin=6pt, innerbottommargin=6pt]{infobox}

\newmdenv[backgroundcolor=bgsuccess, linecolor=success, linewidth=2pt,
leftline=true, rightline=false, topline=false, bottomline=false,
innerleftmargin=10pt, innertopmargin=6pt, innerbottommargin=6pt]{successbox}

\newmdenv[backgroundcolor=bgwarning, linecolor=warning, linewidth=2pt,
leftline=true, rightline=false, topline=false, bottomline=false,
innerleftmargin=10pt, innertopmargin=6pt, innerbottommargin=6pt]{warningbox}

\newmdenv[backgroundcolor=bgdanger, linecolor=danger, linewidth=2pt,
leftline=true, rightline=false, topline=false, bottomline=false,
innerleftmargin=10pt, innertopmargin=6pt, innerbottommargin=6pt]{dangerbox}

\lstset{backgroundcolor=\color{codegray}, basicstyle=\small\ttfamily,
	breaklines=true, frame=single, rulecolor=\color{gray}, numbers=left,
	numberstyle=\tiny\color{gray}, numbersep=5pt, showstringspaces=false}

\newcommand{\rmsleGood}[1]{\textbf{\color{success}#1}}
\newcommand{\rmsleBad}[1]{\textbf{\color{danger}#1}}
\newcommand{\rmsleNeutral}[1]{\textbf{\color{primary}#1}}

% ================================================================
\begin{document}
	\pagenumbering{roman}
	
	\begin{titlepage}
		\centering\vspace*{2cm}
		{\Huge\bfseries\color{primary} Технический отчёт}\\[0.5cm]
		{\Large\bfseries Предсказание LTV пользователей}\\[0.3cm]
		{\large Хакатон по машинному обучению}\\[2cm]
		\rule{\textwidth}{1.5pt}\\[0.4cm]
		{\large\bfseries Задача:} Предсказание суммарного GMV пользователей за 30 дней\\
		на основе обезличенных данных активности в Поиске и Каталоге\\[0.4cm]
		\rule{\textwidth}{0.8pt}\\[2cm]
		\begin{tabular}{ll}
			\textbf{Участник:}    & Кирилл Егоров \\[0.3cm]
			\textbf{Метрика:}     & RMSLE (Root Mean Squared Log Error) \\[0.3cm]
			\textbf{Tie-breaker:} & Gini по предсказаниям + RMSPE суммарного GMV \\[0.3cm]
			\textbf{Данные:}      & 250\,000 пользователей, 01.01.2025--13.02.2026 \\[0.3cm]
			\textbf{Финальный ЛБ:}& \rmsleGood{1.65305} (CB 70\% + MLP 20\% + XGB 10\% $\times 1.12375$) \\[0.3cm]
			\textbf{Публичный ЛБ:}& 20\% (50\,000 клиентов) \\[0.3cm]
			\textbf{Приватный ЛБ:}& 80\% (200\,000 клиентов) \\[0.3cm]
			\textbf{Старт:}       & 10 августа 2026 г. \\[0.3cm]
			\textbf{Дедлайн:}     & $\approx$ 31 августа 2026 г. \\[0.3cm]
			\textbf{Лимит:}       & 5 сабмитов / день, 2 финальных решения \\[0.3cm]
			\textbf{GPU:}         & NVIDIA RTX 4050 Laptop (6 ГБ VRAM) \\
		\end{tabular}
		\vfill
		{\small Дата составления: 28 августа 2026 г. (Финальная редакция)}
	\end{titlepage}
	
	\tableofcontents\newpage
	\pagenumbering{arabic}
	
	% ================================================================
	\section{Описание задачи и данных}
	
	Задача --- предсказание \textbf{Customer Lifetime Value (LTV)} на горизонте 30 дней.
	Особенности целевой переменной: данные сильно \textbf{zero-inflated}. Около 45.1\% пользователей в обучающей выборке не совершают ни одной покупки (target = 0).
	
	\begin{dangerbox}
		\textbf{Публичный / Приватный лидерборд: 20\% / 80\%!}\\
		Оценка на публичном лидерборде вычисляется только по 50\,000 из 250\,000 клиентов.
		Это означает высокий риск \emph{public-private shake-up}. Оптимизация моделей строилась на строгой локальной кросс-валидации (OOF) с последующей ортогональной оптимизацией весов ансамбля и масштабирующих множителей.
	\end{dangerbox}
	
	% ================================================================
	\section{Эволюция признакового пространства и схема валидации}
	Для временных рядов использовался строго \textbf{Walk-Forward CV} ($stride=14$ дней). Схема была расширена до \textbf{5 фолдов} (min\_history=180 дней) для построения устойчивого мета-ансамбля.
	
	\subsection{Исторические версии данных (v1 -- v5)}
	\begin{itemize}[nosep]
		\item \textbf{v1--v2 (Базовые агрегации):} Сбор простых счетчиков действий и сумм. На этом этапе была допущена ошибка "машины времени" (утечка таргета в обучающие признаки), которая привела к ложно-оптимистичным результатам (до 1.69009) и была устранена в v3.
		\item \textbf{v3--v4 (RFM и временные окна):} Внедрение строгой Walk-Forward валидации. Расчет классических RFM-признаков (Recency, Frequency, Monetary) по окнам от 3 до 365 дней.
		\item \textbf{v5 (Взаимосвязи и тренды):} Добавление вторичных признаков, таких как конверсия из поиска в корзину и тренды отношения трат за 7 дней к 30 дням.
	\end{itemize}
	
	\subsection{Финальная версия v6 (Data-Centric AI)}
	Сгенерировано 435 признаков. Ключевые нововведения, обеспечившие максимальный прирост метрики:
	\begin{itemize}[nosep]
		\item \textbf{Межпокупочные паттерны:} mean/std gap между заказами, метрики регулярности (\texttt{order\_regularity\_cv}).
		\item \textbf{Lifecycle Trajectory:} Ускорение заказов (\texttt{order\_acceleration}) и GMV slope.
		\item \textbf{Винзоризация таргета (Target Capping):} Ограничение целевой переменной на этапе обучения по 99.9-му перцентилю (4000) для защиты моделей от аномальных "китов".
		\item \textbf{Dormancy \& Gaps:} Обработка пустых интервалов числом $999.0$ с введением жестких флагов \texttt{is\_single\_buyer} и \texttt{is\_dormant} (спячка > 120 дней).
	\end{itemize}
	
	% ================================================================
	\section{Архитектура Мета-Ансамбля (Stacking)}
	Финальное решение построено на максимизации \textbf{независимости ошибок} (Error Orthogonality). Базовый ансамбль состоит из трех SOTA-моделей. В ходе сканирования (ternary search) весов на чистых данных v6 был подтвержден глобальный локальный оптимум (RMSLE = 1.65559) при распределении:
	
	\subsection{1. CatBoostRegressor (Базовый фундамент, Вес: 70\%)}
	Градиентный бустинг с \emph{level-wise} (oblivious) построением деревьев. Применена техника \textbf{Seed Averaging}: усреднение предсказаний 5 моделей с разными seed-значениями для радикального снижения дисперсии.
	
	\subsection{2. Tabular MLP + ZILN Loss (Вероятностная нейросеть, Вес: 20\%)}
	Полносвязный 4-слойный перцептрон, обученный на сглаженных признаках через \texttt{QuantileTransformer}. Использует функцию потерь ZILN (\emph{Zero-Inflated Lognormal Loss}), которая элегантно и одновременно решает задачи классификации и регрессии в едином пайплайне.
	
	\subsection{3. XGBoost (Охотник на нелинейности, Вес: 10\%)}
	Асимметричный градиентный бустинг с глубокой L1-регуляризацией (\texttt{reg\_alpha = 2.0}).
	
	% ================================================================
	\section{Анализ преодоленных проблем}
	
	\subsection{Взрыв градиентов MLP из-за StandardScaler}
	Применение классического \texttt{StandardScaler} привело к катастрофическому росту Loss. Проблема решена переходом на \textbf{\texttt{QuantileTransformer}} (\texttt{output\_distribution='normal'}), который аккуратно уложил все выбросы в гауссиану.
	
	\subsection{Сезонный сдвиг (Concept Drift) и Праздничные множители}
	Тестовая выборка представляет собой период с 14 февраля по 15 марта (весенние праздники). Модель, обученная на среднегодовых паттернах, систематически недооценивала траты. Аналитика по историческим данным 2025 года показала рост конверсии в праздничный месяц на 3\% (итоговый мультипликатор $\approx 1.09$). 
	
	Применение масштабирующего множителя (до 1.12375) к предсказаниям финального ансамбля позволило пробить потолок чистых моделей и достичь феноменального RMSLE \textbf{1.65305}.
	
	% ================================================================
	\section{Отклоненные гипотезы и стратегии}
	В процессе разработки был протестирован и отвергнут ряд подходов (см. таблицу сабмитов):
	\begin{itemize}
		\item \textbf{Tweedie Loss:} Классическая функция потерь для zero-inflated распределений. Эксперименты с Tweedie в бустингах показали результаты хуже, чем стандартная оптимизация RMSE на логарифмированном таргете $\log(1 + y)$ или использование кастомного ZILN в MLP.
		\item \textbf{Two-Stage Pipeline (Двухэтапная модель):} Идея раздельного обучения бинарного классификатора ($P(y>0)$) и регрессора для покупателей. Была отклонена, так как MLP с ZILN-loss решает эту задачу математически точнее внутри единой архитектуры (совместное обучение $p$ и $\mu$), предотвращая накопление ошибки (error compounding) между моделями.
		\item \textbf{Псевдо-разметка (Pseudo-Labeling):} Эксперимент с добавлением предсказаний с тестовой выборки в тренировочный сет был свернут из-за высокого риска жесткого переобучения на публичный лидерборд (учитывая сплит 20\% / 80\%, доверять паблику на 100\% было опасно).
		\item \textbf{Эмбеддинги Event2Vec (Look-ahead bias):} Обучение Word2Vec на всей истории без строгой Walk-Forward изоляции приводило к заглядыванию в будущее. Честная изоляция эмбеддингов внутри фолдов не дала ожидаемого буста метрики из-за высокой разреженности последовательностей покупок.
		\item \textbf{Вероятностные модели BTYD (Buy Till You Die):} Модели вроде BG/NBD показали крайне слабую способность предсказывать LTV на коротком горизонте (30 дней) для наших данных (RMSLE > 1.76).
		\item \textbf{Логарифмический блендинг:} Усреднение логарифмов предсказаний (геометрическое среднее) слишком сильно занижало прогнозы для крупных покупателей по сравнению с классическим арифметическим усреднением.
	\end{itemize}
	
	% ================================================================
	\section{Полная история сабмитов (Leaderboard)}
	Ниже задокументированы все 68 отправленных решений, отсортированные от лучшей метрики к худшей. 
	
	\begin{table}[h!]
		\centering\small
		\caption{Часть 1: Чемпионы и финальные тюнинги (Топ-24)}
		\vspace{0.2cm}
		\begin{tabular}{l l p{7cm}}
			\toprule
			\textbf{Имя файла} & \textbf{RMSLE} & \textbf{Примечание} \\
			\midrule
			\texttt{v6\_champion\_mult\_1.12375.csv} & \rmsleGood{1.65305} & \textbf{Абсолютный оптимум (Множитель 1.12375)} \\
			\texttt{v6\_champion\_mult\_1.125.csv} & \rmsleGood{1.65305} & Бинарный поиск множителя \\
			\texttt{v6\_champion\_mult\_1.1225.csv} & \rmsleGood{1.65305} & Бинарный поиск множителя \\
			\texttt{v6\_champion\_mult\_1.12.csv} & \rmsleGood{1.65306} & Разведка правой границы \\
			\texttt{v6\_champion\_mult\_1.13.csv} & \rmsleGood{1.65306} & Разведка правой границы \\
			\texttt{v6\_champion\_mult\_1.14.csv} & \rmsleGood{1.65309} & Подтверждение экстремума \\
			\texttt{v6\_champion\_mult\_1.10.csv} & \rmsleGood{1.65313} & Приближение к дну \\
			\texttt{v6\_champion\_mult\_1.09.csv} & \rmsleGood{1.65321} & Расчетный сдвиг 2025 года \\
			\texttt{v6\_champion\_mult\_1.05.csv} & \rmsleGood{1.65389} & Успешный старт калибровки \\
			\texttt{v6\_champion\_mult\_1.03.csv} & \rmsleGood{1.65445} & \\
			\texttt{v6\_champion\_mult\_1.02.csv} & \rmsleGood{1.65479} & \\
			\texttt{v6\_arithmetic\_cb70\_mlp20\_xgb10} & \rmsleNeutral{1.65559} & \textbf{Лучший чистый базовый ансамбль} \\
			\texttt{v6\_arithmetic\_cb65\_mlp25\_xgb10} & 1.65566 & Тест увеличенной доли MLP \\
			\texttt{v6\_arithmetic\_cb65\_mlp20\_xgb15} & 1.65568 & Тест увеличенной доли XGBoost \\
			\texttt{v6\_arithmetic\_cb75\_mlp15\_xgb10} & 1.65574 & \\
			\texttt{arithmetic\_cb75\_mlp15\_xgb10.csv} & 1.65581 & \\
			\texttt{arithmetic\_cb75\_mlp20\_xgb05.csv} & 1.65584 & \\
			\texttt{FINAL\_blend\_cb80\_mlp15\_xgb05} & 1.65590 & \\
			\texttt{arithmetic\_cb85\_mlp10\_xgb05.csv} & 1.65601 & \\
			\texttt{FINAL\_blend\_cb85\_mlp15.csv} & 1.65605 & \\
			\texttt{v6\_champion\_mult\_0.98.csv} & 1.65655 & Ухудшение (занижение прогноза) \\
			\texttt{catboost\_v5\_seed\_avg.csv} & 1.65656 & \\
			\texttt{catboost\_v6\_seed\_avg.csv} & 1.65678 & Соло CatBoost (v6) \\
			\texttt{v6\_champion\_mult\_0.97.csv} & 1.65710 & \\
			\bottomrule
		\end{tabular}
	\end{table}
	
	\begin{table}[h!]
		\centering\small
		\caption{Часть 2: Поиск базового ансамбля и эксперименты (25–46)}
		\vspace{0.2cm}
		\begin{tabular}{l l p{7cm}}
			\toprule
			\textbf{Имя файла} & \textbf{RMSLE} & \textbf{Примечание} \\
			\midrule
			\texttt{blend\_seed\_avg\_92\_rnn8.csv} & 1.65878 & Бывший лидер (v4b + RNN) \\
			\texttt{blend\_seed\_avg\_90\_rnn10.csv} & 1.65884 & \\
			\texttt{blend\_seed\_avg\_95\_rnn5.csv} & 1.65890 & \\
			\texttt{frank\_cb70\_ziln20\_xgb5\_rnn5.csv} & 1.65912 & \\
			\texttt{frank\_cb75\_ziln20\_rnn5.csv} & 1.65912 & \\
			\texttt{blend\_pseudo\_90\_rnn10.csv} & 1.65929 & Эксперимент с pseudo-labeling \\
			\texttt{blend\_ortho\_90\_10.csv} & 1.65935 & \\
			\texttt{logblend\_cb800\_mlp150\_xgb50.csv} & 1.65975 & Отказ от Log-Space блендинга \\
			\texttt{frank\_cb70\_ziln15\_rnn15.csv} & 1.65977 & \\
			\texttt{catboost\_v4a\_seed\_avg.csv} & 1.65984 & \\
			\texttt{blend\_ortho\_85\_15.csv} & 1.65992 & \\
			\texttt{ziln\_v6\_v1\_submission.csv} & 1.66008 & Рекорд соло-нейросети ZILN (v6) \\
			\texttt{catboost\_v4a\_baseline\_sub...} & 1.66049 & \\
			\texttt{blend\_v4a\_ortho\_95\_05.csv} & 1.66071 & \\
			\texttt{blend\_v4a\_ortho\_90\_10.csv} & 1.66113 & \\
			\texttt{ziln\_v5\_v1\_submission.csv} & 1.66220 & \\
			\texttt{catboost\_v5\_seed\_avg.csv} & 1.66265 & Оверфит Optuna \\
			\texttt{blend\_v4a60\_lgbm40.csv} & 1.66336 & \\
			\texttt{xgboost\_v6\_submission.csv} & 1.66478 & Соло XGBoost (v6) \\
			\texttt{xgboost\_v5\_submission.csv} & 1.66506 & \\
			\texttt{mlp\_ziln\_submission.csv} & 1.66545 & \\
			\texttt{catboost\_baseline\_submission} & 1.67028 & Baseline от организаторов \\
			\bottomrule
		\end{tabular}
	\end{table}
	
	\begin{table}[h!]
		\centering\small
		\caption{Часть 3: Ранние версии, утечки и бейзлайны (47–68)}
		\vspace{0.2cm}
		\begin{tabular}{l l p{7cm}}
			\toprule
			\textbf{Имя файла} & \textbf{RMSLE} & \textbf{Примечание} \\
			\midrule
			\texttt{xgboost\_log\_submission.csv} & 1.67102 & \\
			\texttt{lgbm\_v3\_submission.csv} & 1.67115 & \\
			\texttt{catboost\_v3\_submission.csv} & 1.67116 & \\
			\texttt{catboost\_v4\_cumul\_sub...} & 1.67126 & \\
			\texttt{catboost\_v4\_single\_sub...} & 1.67177 & \\
			\texttt{blend\_cbv3\_80\_rnnv3\_19.csv} & 1.67380 & \\
			\texttt{blend\_cbv3\_70\_rnnv3\_30.csv} & 1.67583 & \\
			\texttt{ensemble\_cb70\_rnn30.csv} & 1.67588 & \\
			\texttt{blend\_cbv3\_65\_rnnv3\_35.csv} & 1.67702 & \\
			\texttt{lgbm\_v4\_submission.csv} & 1.68433 & Нестабильность LightGBM \\
			\texttt{ensemble\_v2\_cb70\_rnn30.csv} & 1.68803 & Утечка данных в v2 \\
			\texttt{ensemble\_v2\_cb60\_rnn40.csv} & 1.68942 & Утечка данных в v2 \\
			\texttt{catboost\_v2\_submission.csv} & 1.69009 & Утечка данных в v2 \\
			\texttt{lgbm\_v4\_submission.csv} & 1.70118 & \\
			\texttt{rnn\_v4\_submission.csv} & 1.70489 & Соло RNN (v4) \\
			\texttt{rnn\_v2\_submission.csv} & 1.71792 & \\
			\texttt{catboost\_v4a\_zero\_hacked} & 1.73573 & Провал гипотезы жестких нулей \\
			\texttt{rnn\_baseline\_submission.csv} & 1.73624 & \\
			\texttt{btyd\_blend\_btyd2\_cb8.csv} & 1.76655 & Провал вероятностной модели BTYD \\
			\texttt{catboost\_tuned\_submission} & 1.79400 & Провал поиска параметров Optuna v1 \\
			\texttt{btyd\_blend\_btyd3\_cb7.csv} & 1.81827 & \\
			\texttt{rnn\_v4\_submission.csv} & \rmsleBad{1.87947} & Критическая ошибка/утечка таргета \\
			\bottomrule
		\end{tabular}
	\end{table}
	
\end{document}